\section{Related work}\label{sec:related}
\subsection{Dynamic task allocation and routing}
Multi-robot task allocation~\cite{gerkey2004formal} and dynamic vehicle
routing~\cite{pillac2013review,psaraftis2016dynamic} reassign tasks when information changes. The
planning problem considered here is a prize-collecting routing problem with time
windows~\cite{solomon1987algorithms,vansteenwegen2011orienteering}. Large neighbourhood search has been adapted to dynamic routing with time windows~\cite{alnsdvrp2024}. For multi-UAV
systems, auction methods such as CBBA reach conflict-free assignments quickly~\cite{choi2009consensus}.
Recent work handles stochastic task completion with hierarchical policies~\cite{choudhury2022dynamic}, uncertain
execution times with dynamic time slices~\cite{huang2024dynamic}, intermittent communication with uncertainty-gated
distributed allocation~\cite{ucdg2026drones}, and events with assignment-preserving or local
repair~\cite{hwang2026apr,shen2026localrepair}. Event-triggered and periodic allocation have been compared for
multi-UAV systems~\cite{cong2026event}, in line with the rescheduling
literature~\cite{vieira2003rescheduling,ouelhadj2009survey}. Calvo and Capitan validate a repair-based architecture
with PX4 and Gazebo~\cite{calvo2025heterogeneous}, and Ali et al.\ use ArduPilot software-in-the-loop
simulation~\cite{ali2026sitl}. Table~\ref{tab:related} compares these approaches.

\inputtable{tab_related}

\subsection{Uncertainty in routing and allocation}
Stochastic vehicle routing models travel and service times as random
variables~\cite{gendreau1996stochastic}. With hard time windows and stochastic service times, a priori routes with
recourse~\cite{errico2016apriori} and robust formulations over budgeted uncertainty
sets~\cite{agra2013robust} protect the schedule against delays, at the price of a harder model. Scenario-based
approaches plan against sampled futures~\cite{bent2004scenario}, and chance-constrained task allocation bounds the
probability that a team misses its objectives~\cite{ponda2012chance}. Per-customer chance constraints on the service
time windows were formulated for routing with stochastic travel and service times by Li et al.~\cite{li2010vrptw}.
Ehmke et al.~\cite{ehmke2015ensuring} propagate the mean and variance of the arrival time along a route, approximate
the waiting for a time window by Clark's moments of a maximum~\cite{clark1961greatest}, and use the result as a
feasibility check inside routing heuristics. Soft time windows take the opposite route and charge a penalty for every
late service inside the objective. Ta\c{s} et al.~\cite{tas2013stochastic} minimise transportation costs plus expected
early and late service costs under stochastic travel times with tabu search, which requires the distribution of the
lateness of every route during the search. For multi-UAV systems, Gao et al.~\cite{gao2025drcc} formulate a
distributionally robust chance-constrained task assignment with time windows under uncertain fuel consumption and
solve a mixed-integer reformulation, Wang et al.~\cite{wang2024uncertain} estimate the effect of uncertain flight times
on completion times and time-window violations without Monte Carlo simulation, and Zhou and
Tokekar~\cite{zhou2022risk} select risk-averse multi-robot assignments through the conditional value at risk. These
formulations rely on mixed-integer, sampling-based or greedy solvers over a planning horizon, and none of them is
evaluated as a replanner that has to publish a plan within a fraction of a second after every event.

The proposed model differs in four respects. Its error model is fitted to the finish times observed in closed loop
rather than assumed per arc. It is applied in a replanner that runs between events within a fraction of a second, with
a publication delay. It adds a fallback for the tasks that a hard chance constraint gives up. And it prices a late
arrival only in this fallback, through a break-even probability that again yields a deterministic deadline, so the
search keeps its objective and its fast feasibility checks, whereas soft time windows change both. Moment propagation,
a proportional buffer, a constant buffer and the consensus-based bundle algorithm (CBBA)~\cite{choi2009consensus} are
re-implemented as baselines.

\subsection{Metareasoning and anytime algorithms}
Anytime algorithms~\cite{dean1988analysis,zilberstein1996using} and the rational allocation of
computation~\cite{russell1991right,boddy1994deliberation} decide how long a planner should compute, from performance
profiles~\cite{hansen2001monitoring}, online performance prediction~\cite{svegliato2018meta} or learned stopping
policies~\cite{sung2021quit}. Situated temporal planning and deadline-aware search account for the passage of time
during planning~\cite{cashmore2018clock,shperberg2021situated,dionne2011deadline}. Two of these rules are
re-implemented as baselines, so that the value of controlling computation can be compared with the value of the
planning model.
